\clearpage
\VerdierPhysicalPageAnchor{0132}%
\begingroup
\small
\setlength{\emergencystretch}{1.5em}%
\setlength{\abovedisplayskip}{2pt plus 1pt minus 1pt}%
\setlength{\belowdisplayskip}{2pt plus 1pt minus 1pt}%

\VerdierSemanticAnchor{ch2-2-1-10-continuation}
\noindent suite, SM3).
\VerdierSemanticAnchor{ch2-2-1-10-triangle-setup}%
Par hypothèse, il existe un triangle distingué
$(X,Y,N,s,h,u)$, où $N$ est un objet de $\mathcal D'$. Comme on a $fs=0$, il
résulte de (1.2.1) qu'il existe un morphisme $g\colon N\to Z$ tel que $f=gh$.

\VerdierSemanticAnchor{ch2-2-1-10-factorization}
Soit alors $(N,Z,W,g,t,v)$ un triangle distingué contenant $g$ (TRI). On a
alors $tg=0$ (1.2.2) et, par suite, $tf=tgh=0$.

\VerdierSemanticAnchor{ch2-2-1-10-closure-conclusion}
Tout revient donc à montrer que $t$ est un élément de $S(\mathcal D')$, ce
qui résulte immédiatement de TRII).

\VerdierSemanticAnchor{ch2-2-1-11}
\medskip
\noindent\textbf{2.1.11.}\quad
\VerdierSemanticAnchor{ch2-2-1-11-sm5}
\emph{L'ensemble $S(\mathcal D')$ possède la propriété} SM5).
\VerdierSemanticAnchor{ch2-2-1-11-proof}
Résulte immédiatement de TRII) et du fait que la sous-catégorie $\mathcal D'$
est stable par translation.

\VerdierSemanticAnchor{ch2-2-1-12}
\medskip
\noindent\textbf{2.1.12.}\quad
\VerdierSemanticAnchor{ch2-2-1-12-sm6}
\emph{L'ensemble $S(\mathcal D')$ possède la propriété} SM6).
\VerdierSemanticAnchor{ch2-2-1-12-triangle-pair}
Soient $(X,Y,Z,u,v,w)$ et $(X',Y',Z',u',v',w')$ deux triangles distingués de
$\mathcal D$ et :

\VerdierSemanticAnchor{ch2-2-1-12-input-square}
\begin{center}
  \begin{tikzcd}[column sep=4.2em,row sep=2.5em]
    X \arrow[r,"u"] \arrow[d,"s"'] &
    Y \arrow[d,"{s'}"] \\
    X' \arrow[r,"{u'}"] & Y'
  \end{tikzcd}
\end{center}

\VerdierSemanticAnchor{ch2-2-1-12-completion-goal}
\noindent un diagramme commutatif, où $s$ et $s'$ appartiennent à
$S(\mathcal D')$. Il s'agit de trouver un morphisme $s''\colon Z\to Z'$ de
$S(\mathcal D')$ tel que $(s,s',s'')$ soit un morphisme de triangles.

\VerdierSemanticAnchor{ch2-2-1-12-case-reduction}
On peut tout d'abord supposer qu'on est dans l'un des deux cas suivants :

\begin{enumerate}
\renewcommand{\labelenumi}{\emph{\alph{enumi})}}
  \item \VerdierSemanticAnchor{ch2-2-1-12-case-a}%
  On a $X=X'$ et $s=\operatorname{id}_X$.
  \item \VerdierSemanticAnchor{ch2-2-1-12-case-b}%
  On a $Y=Y'$ et $s'=\operatorname{id}_Y$.
\end{enumerate}

\VerdierSemanticAnchor{ch2-2-1-12-factorization}
En effet, tout morphisme de morphismes se factorise :

\VerdierSemanticAnchor{ch2-2-1-12-factorization-diagram}
\begin{center}
  \begin{tikzcd}[column sep=4.2em,row sep=1.7em]
    X \arrow[r,"u"] \arrow[d,"\operatorname{id}_X"'] &
    Y \arrow[d,"{s'}"] \\
    X \arrow[r,"{s'u}"] \arrow[d,"s"'] &
    Y' \arrow[d,"\operatorname{id}_{Y'}"] \\
    X' \arrow[r,"{u'}"] & Y'
  \end{tikzcd}
  \quad ,
\end{center}

\VerdierSemanticAnchor{ch2-2-1-12-composition}
\noindent et on obtiendra le morphisme $s''$ en composant (on a déjà montré
que l'ensemble $S(\mathcal D')$ est stable par composition).

\VerdierSemanticAnchor{ch2-2-1-12-opposite-category}
Ensuite, le cas b) se déduit du cas a) en passant aux catégories opposées.

\VerdierSemanticAnchor{ch2-2-1-12-triv-conclusion}
Il suffit donc de traiter le cas a), qui est alors une conséquence immédiate
de TRIV).

\VerdierSemanticAnchor{ch2-2-1-13-p133-boundary}
\endgroup
